\section{Related Work}
\label{sec:related}

\subsection{Logical and physical time}

Lamport clocks and happened-before relations establish that causal order is not
the same as physical timestamp order \cite{lamport1978time}. Vector-clock
systems make concurrent events explicit \cite{mattern1989virtual}. The
observer-declared timeline rule does not replace these mechanisms. It brings a
similar discipline to product surfaces: when the product shows a timeline, it
should expose enough perspective metadata for the ordering to be audited.

Distributed snapshots construct a consistent global state without requiring a
physically simultaneous observation \cite{chandy1985snapshots}. An
observer-declared view has a related concern, but its contract extends beyond
one snapshot algorithm: it names accepted heterogeneous facts, projection
semantics, evidence degradation, and the incarnation that claims the view.

\subsection{Failure detection and recovery epochs}

Failure detectors reason about what distributed processes can suspect and what
reliable services can be built from those suspicions \cite{chandra1996failure}.
The incarnation model addresses a product-layer consequence: replacement must
not turn liveness evidence into authority continuity. Generation and epoch
fences are implementation techniques for rejecting stale actors; the proposed
surface additionally asks a timeline to expose the cut and semantic relation
that make a successor view reproducible.

\subsection{Event sourcing}

Event sourcing treats state as a derivation from an append-only sequence of
events \cite{fowler2005eventsourcing}. Observer-declared timelines share the
idea that records matter, but differ in emphasis. The visible timeline is not
assumed to be the authoritative event log. It is a perspective-bearing
projection over accepted facts that may come from several sources. Episodes add
a bounded causal and verification object above physical append blocks, while
keeping projections rebuildable from authority.

\subsection{Distributed tracing}

Distributed tracing systems such as OpenTelemetry capture spans and causal
relationships across services \cite{opentelemetry}. This is valuable evidence
for runtime behavior. The proposed rule is broader at the product surface: a
timeline may combine traces, repository facts, user decisions, release
manifests, local files, and remote agent actions. An Episode may contain trace
evidence, but its boundary is defined by causal closure and portable trust rather
than by one tracing backend or span vocabulary.

\subsection{CRDTs and replicated state}

CRDTs provide convergence properties for replicated data types under
concurrent updates \cite{shapiro2011crdt}. Observer-declared timelines are not
a conflict-free data structure. They are a view contract: even when underlying
state uses CRDTs, consensus, or a central sequencer, a mixed-source product
timeline should still declare what perspective the user or agent is seeing.

\subsection{Human-computer and human-agent cooperation}

Engelbart framed computing as augmenting human intellect rather than merely
automating tasks \cite{engelbart1962augmenting}. Real-world agent work extends
that problem. The product must support cooperation between human and agent
participants that observe, remember, and act through different surfaces.
